\documentclass[../../main.tex]{subfiles}
\begin{document}
\begin{oframed}
  \subfile{problem}
\end{oframed}

\begin{flushright}
  \footnotesize\textit{（(a)(b)いずれか一方を選択する問題．(a)のみ解答あり．(b)は原稿は空欄）}
\end{flushright}

{\bf [解答]}
(a)-1

まず最初に任意の $n\in\mathbb{N}$ に対して$|a_n|\le1$ となることを数学的帰納法で示す．
$a_N=1$ 及び $a_{n+1}=2a_n^2-1$ だから $a_{N+1}=a_{N+2}=\cdots = 1$であり，$n\ge N$ の時，$a_n=1$ で成立する．

そこで以下 $n\le N$ の時を考える．$n=k$（$2\le k\le N$）での成立を仮定すると漸化式より
\begin{align}
  a_{k-1}^2=\dfrac{1}{2}(a_k+1)
  \therefore\
  0\le a_{k-1}^2\le 1 \quad (\because|a_k|\le 1)
  \therefore\
  |a_{k-1}|\le 1
\end{align}
となり $n=k-1$ でも成立する．
したがって帰納法により，任意の $n\in\mathbb{N}$ に対して $|a_n|\le 1$ となる．

よって$n=1$とすれば題意の主張 $|a_1|\le 1$ が従う．

(a)-2
(a)-1より$|a_n|\le 1$だから
\begin{align}
  a_n=\cos\theta_n
\end{align}
とおける．漸化式に代入すると倍角公式により
\begin{align}
  \cos\theta_{n+1} &=\cos2\theta_n
\end{align}
だから繰り返し用いて
\begin{align}
  \cos\theta_{n} = \cos \left(2^{n-1}\theta_1\right) \label{eq:1}
\end{align}
と書ける．% TODO: ここはもう少し丁寧に書いても良いかも

ここで$a_N=1$より\cref{eq:1}に$n=N$を代入すると
\begin{align}
  1 = \cos \left(2^{N-1}\theta_1\right)
  \quad \therefore
  \theta_1=\dfrac{2k\pi}{2^{N-1}}=\dfrac{k\pi}{2^{N-2}}
  \quad (k \in\mathbb{Z})
\end{align}
と書ける．よって題意は示された．

\end{document}
